\documentclass{article}
\usepackage{tikz}
\usetikzlibrary{angles, quotes}
\usepackage{amssymb}
\usepackage{pgfplots}
\pgfplotsset{width=10cm,compat=1.9}
\usepgfplotslibrary{statistics}
\usepackage[left=2.54cm,right=2.54cm,top=2.54cm,bottom=2.54cm]{geometry}
\usepackage{tasks} %Paquete necesario para la producción de items horizontales para algunos ejercicios de tarea empleando el comando /begin{multicols}}{#}.../end{multicols} para ayudarnos a reducir las páginas
\usepackage{pdflscape} %comando que permite cambiar a orientación horizontal las hojas%
\usepackage{soul} %Paquete para permitir la compilación de acentos usando el comando {\hl{}}}\\
\sethlcolor{yellow(munsell)} %comando para definir el color para el subrayado usando el comando \hl
\usepackage[utf8]{inputenc}
\usepackage[spanish]{babel}
\usepackage{icomma}
\usepackage{ragged2e}
\usepackage{siunitx}
\usepackage{url}
\usepackage[colorlinks=true, urlcolor=blue,  linkcolor=black, citecolor=green]{hyperref}
\usepackage{pdfpages}
\usepackage{blindtext} %Paquete que genera texto ficticio (dummy text) con el comando: \blindtext
\usepackage{cancel} %in the preamble gives you four different modes of striking through: \cancel{text to cancel} draws a diagonal line (slash) through its argument, \bcancel{text to cancel} uses the negative slope (a backslash), \xcancel{text to cancel} draws an X (actually \cancel plus \bcancel), \cancelto{〈value〉}{〈expression〉} draws a diagonal arrow through the 〈expression〉pointing to the 〈value〉 (math-mode only)
\usepackage{csquotes}
\usepackage{afterpage}
\usepackage{parskip} 
\usepackage{float}
\usepackage{enumitem}
\usepackage{multicol}%Paquete que permite la creación aislada de columnas de texto. De esta forma se reduce la cantidad de páginas en nuestro documento
\usepackage{lipsum} %Paquete que genera texto de relleno al igual que \usepackage{blindtext}, se usa con el comando \lipsum[#-#] los #(hashtags) delimitan el número de páginas que deseamos ocupar del paquete lipsum para usarlos como relleno. 

% Definir el comando personalizado para la nota
\newcommand{\mynote}[1]{%
\begin{tikzpicture}[baseline=-0.75ex]
\node[draw, circle, fill=Apple Green!60, inner sep=2pt] (note) {\textbf{Nota:}};
\end{tikzpicture}%
\ \textit{#1}%
}

\newenvironment{Figura}
{\par\medskip\noindent\minipage{\linewidth}}
{\endminipage\par\medskip}
\usepackage{caption}
\usepackage[
backend=biber,
sorting=none,
url=true
]{biblatex} %bibliografía
\addbibresource{biblio.bib}
% Establecer el espacio entre las entradas de la bibliografía
\setlength{\bibitemsep}{1\baselineskip} % Puedes ajustar el valor según tus preferencias
\usepackage{amsmath,amsthm,amssymb,amsfonts}
\usepackage{pifont} %Permite la compilación del comando \xmark: es el símbolo de la tachita
\usepackage{mathtools} %permite la compilación de símbolos para matrices
\usepackage{empheq} %Paquete que se relaciona con \usepackage[most]{tcolorbox} para la creación de cuadros/cajas de colores para delimitar los resultados matemáticos o líneas de texto usando la línea de comando: \begin{empheq}[box={\mymath[colback=orange(webcolor)!70,drop lifted shadow]}]{equation*} \end{empheq}
\usepackage[most]{tcolorbox} %Paquete que se relaciona con \usepackage{empheq} para la creación de cuadros/cajas de colores para delimitar los resultados matemáticos o líneas de texto
\renewcommand{\qedsymbol}{$\blacksquare$}
\usetikzlibrary{positioning,decorations.pathreplacing} %Paquete que permite la compilación de llaves para cuadros sinópticos (brace diagram) 
\usepackage{schemata} %The schemata package is designed just for these brace diagrams
\usepackage{booktabs}

\newcommand\AB[2]{\schema{\schemabox{#1}}{\schemabox{#2}}} %comando o paquete necesario para crear cuadros sinópticos

\newtcbox{\mymath}[1][]{%
nobeforeafter, math upper, tcbox raise base,
enhanced, colframe=blue!30!black,
colback=blue!30, boxrule=1pt,
#1} %Comando importante para encasillar los resultados matemáticos en cajas de diferentes colores, se relaciona con los paquetes \usepackage{empheq} y \usepackage[most]{tcolorbox} 

\newenvironment{sysmatrix}[1]
{\left(\begin{array}{@{}#1@{}}}
{\end{array}\right)}
\newcommand{\ro}[1]{%
\xrightarrow{\mathmakebox[\rowidth]{#1}}%
}
\newlength{\rowidth}% row operation width
\AtBeginDocument{\setlength{\rowidth}{3em}}  %Comando importante para laproducción de lineas de operaciones entre matrices, método de Gauss o Gauss Jordan se relaciona con \newenvironment{sysmatrix}

%formato para cambiar el horario a español
\usepackage[spanish]{datetime2}
\DTMsetdatestyle{spanish}

\renewcommand{\today}{\DTMdisplaydate{\the\year}{\the\month}{\the\day }{-1}}
%%%%%%%%%%%%%%%%%%%%%%%%%%%%%%%%%%%%%%%%%

\usepackage{datetime}
\newcommand{\mycurrenttime}{\xxivtime}

%%%%%%%%%%%%%%%%%%%%%%%%%%%%%%%%%%%%%%%%%%%%%%%%%%%%%%%%%%%%%%%%%%%%%%%%%%%%%%%%%%%%%%%%%%%%%%%%%%%%%%%%%%%%%%%%%%%%%%%%%%%%%%%%%%%%%%%%%%%%%%%%%


%%%%%%%%%%%%%Caja de color para el título%%%%%%%%%%%%%%%%%%%%%%%%%%%%%%%%%%%%%%%%%%%%%%%%%%%%%%

\definecolor{myframecolor}{RGB}{1, 45, 75} % Prussian Blue
\definecolor{myboxcolor}{RGB}{0, 107, 92} % Apple Green

\newtcolorbox{mybox}{
enhanced,
colback=myboxcolor!7, % Color de fondo del cuadro
colframe=myframecolor, % Color del marco del cuadro
arc=0pt,
boxrule=1pt,
borderline west={2mm}{-10mm}{myframecolor}, % Borde en el lado izquierdo
sharp corners=southwest,
width=\linewidth,
before=\par\vspace{\bigskipamount}, % Espacio antes del cuadro
after=\par\vspace{\bigskipamount} % Espacio después del cuadro
}


\newcommand{\euler}{e} %Comando para producir letra e de euler
\newenvironment{remark}{\par\vfill\footnotesize % Vertical white space above the remark and smaller font size
\begin{list}{}{
	\leftmargin=80pt % Indentation on the left
	\rightmargin=60pt}\item\ignorespaces % Indentation on the right
\makebox[-2.5pt]{\begin{tikzpicture}[overlay]
		\node[draw=Horizon!60,line width=2.5pt,circle,fill=Horizon!25,font=\sffamily\bfseries,inner sep=4pt,outer sep=0pt] at (-19pt,5pt){\textcolor{Horizon}{Nota}};\end{tikzpicture}} % Blue Nota in a circle
\advance\baselineskip -1pt}{\end{list}\vskip5pt} % Tighter line spacing and white space after remark
\usepackage{graphicx}

\usepackage{titling}

\input{colores.tex}

%%Producir signos de cita textual``Asignatura''

\newcommand{\xmark}{\ding{55}} %%%comando que genera la tachita

\newenvironment{MyColorPar}[1]{%
\leavevmode\color{#1}\ignorespaces%
}{%
}%

%%%%%%%%%%%%%%%%%%% Título de la tarea, Nombre de alumno

\title{Clase \# 35 - Mecánica Estadística}
\author{\emph{JA} $\boldsymbol{\mid}$}
\date{\today}

%%%%%%%%%%%%%%%%%%% Datos de la Materia

\usepackage{fancyhdr}
\fancypagestyle{plain}{%  the preset of fancyhdr 
\fancyhf{} % clear all header and footer fields
\fancyfoot[R]{
\includegraphics[width=2cm]{LogoFCFMBUAP (1).png}}
\fancyfoot[L]{{\bfseries{\thedate{}}} a las {\bfseries{\mycurrenttime{}}} horas{} (GMT-6, H. Puebla de Zaragoza, Pue)}
\fancyhead[L]{Mecánica Estadística}
\fancyhead[R]{\theauthor}
}
\makeatletter
\def\@maketitle{%
\newpage
\null
\vskip 1em%
\begin{center}%
\let \footnote \thanks
{\LARGE \@title \par}%
\vskip 1em%
%{\large \@date}%
\end{center}%
\par
\vskip 1em}
\makeatother

\usepackage{lipsum}  
\usepackage{cmbright}


%%%%%%%%%%%%%%%%%%%%%%%%%%%%%%%%%%%%%%%%%%%%%%%%%%%%%%%%%%%%%%%%%%%%%%%%%%%%%%%%%%%%%%%%%%%%%%%%%%%%%%%%%%%%%%%%%%%%%%%%%%%%%%%%%%%%%%%%%%%%%%%%%%%%%%%%%%%%%%%%%%%%%%%%%%%%%%%%%%%%%%%%%%%%%%%%%%%%%%

\begin{document}



\maketitle


%%%%%%%%%%%%%%%%%%% Datos del profesor, campus y aula

\begin{mybox}
\noindent\begin{tabular}{@{}ll}
	{\bfseries{Profesor}} & Dr. Jorge Velázquez Castro\\
	\textcolor{prussianblue}{\bfseries{Campus}} / \textcolor{trueblue}{\bfseries{Facultad}}  & \textcolor{prussianblue}{\bfseries{C.U. BUAP}} /  \textcolor{trueblue}{\bfseries{FCFM}} \\
	\textcolor{prussianblue}{\bfseries{Edificio}} /  \textcolor{trueblue}{\bfseries{Salón}}    & 	\textcolor{prussianblue}{\bfseries{EMA4}} /  \textcolor{trueblue}{\bfseries{407}}   
\end{tabular} 
\end{mybox} \vspace{1cm}

\tableofcontents

\vspace{2cm}

Anteriormente

Podemos encontrar el primedio de la componete $z$ del momento magnético:

Tenemos:


\[ \langle \mu_{z}  \rangle = \frac{1}{\beta}  \frac{\partial }{\partial H} \ln z (T,H,1) \]

Esto es:

\[   = - \frac{\partial}{\partial H} F(T,H,1)   \]

\fbox{Interesante}


El momento dipolar total


\[ \langle D_{z} \rangle = - \frac{\partial}{\partial H} F(T,H,N) \]

Con:


\[ F(T,H,N) = -N K_{B}\ln \left[    4\pi \frac{\sen \beta \mu H}{4\pi \sen \beta \mu H}    \right]  \]

Tenemos:

\[  \underbrace{\mu_{z}}_{extensiva} \quad \text{es la variable conjugada a} \quad \underbrace{H}_{intensiva} \]

\section{De forma explicícita}

Tenemos:

\[ \langle  D_{z} \rangle = N \langle \mu_{z} \rangle = N \mu \left[ \underbrace{\cot{h} (\beta \mu H) -  \frac{1}{\beta \mu H}}_{F.  de Langevin} \right] \]




\begin{figure}[H]
    \centering
    \includegraphics[width=0.7\textwidth]{1.png}
    \caption{.}
    \label{fig:derivada_total_H}
\end{figure}


SI se hace una expansión en serie de taylor de:


\[ L(x) \approx \frac{x}{3} - \frac{x^{3}}{45} + \ldots  + \]


Es decir, a tmperaturas altas:

\[   x<<1 \]

Tenemos:

\[  \langle  D_{z}\rangle \approx \left. \frac{N\mu^{2}}{3K_{B}T} H \right\} \text{Ley de Curie} \]


Con:

\[\boxed{ \color{trueblue} D_{z} = \frac{C}{T}H  }  \]


La \textbf{susceptibilidad magnética} es:



\[ \mathcal{X} = \lim\limits_{H \to 0} \frac{  \partial \langle D_{z} \rangle  }{\partial H } = \frac{N  \mu^{2} }{3K_{B}T} = \overbrace{\frac{C}{T}}^{cte. de curie} \]

\[ C = \frac{N \mu^{3}}{3K_{B}}  \]


\subsection{Paramagnetismo cuántico}

En el tratamiento cuántico del paramagnetismo, la energía de interacción entre un momento magnético y un campo externo uniforme \(\mathbf{H}\) está dada por
\[
E_{m} = -\,g\,\mu_{B}\,H\,m,
\]
donde \(\mu_{B}\) es el \textbf{magnetón de Bohr}, definido como
\[
\mu_{B} = \frac{e\hbar}{2mc},
\]
y \(g\) es el \textbf{factor giromagnético} característico de la partícula.

El número cuántico \(m\) toma los valores
\[
m = -j,\, -j+1,\, \ldots,\, j-1,\, j,
\]
donde \(j\) representa el \textbf{momento angular total} (o espín) de la partícula.

\textbf{Interpretación física.}  
Cada partícula con momento angular \(j\) puede orientarse en \(2j+1\) direcciones espaciales distintas con respecto al campo magnético.  
Por ejemplo:
\begin{itemize}
  \item Para el \textbf{electrón} (\(j = \tfrac{1}{2}\)), existen dos orientaciones posibles: \(m = \pm \tfrac{1}{2}\).
  \item Para una partícula de \textbf{espín 1}, hay tres orientaciones: \(m = -1,\,0,\,1\).
  \item Para \(j = \tfrac{3}{2}\), existen cuatro orientaciones: \(m = -\tfrac{3}{2},\,-\tfrac{1}{2},\,\tfrac{1}{2},\,\tfrac{3}{2}\).
\end{itemize}

Estas distintas orientaciones corresponden a \textbf{niveles de energía discretos}, separados por intervalos proporcionales al campo magnético aplicado.  
Así, el campo \(\mathbf{H}\) rompe la degeneración de los estados magnéticos, fenómeno conocido como \textbf{efecto Zeeman}.  
En el contexto del paramagnetismo, esta cuantización del momento magnético permite obtener la magnetización media mediante la estadística de Boltzmann y conduce directamente a la \textbf{ley de Curie} en el límite clásico.


\textit{Una cosa es el spin y otra cosa es el momento angular intrinseco. Si tenemos la propiedad de que la partucka rote entonces la partícula va a generar un campo magnético, entonces debe de existir una relación entre esas cantidades. Debe de existir una cantidad que mapee el spin con el momento intrinseco y precisamente es el factor giromagnético $g$.}

\textit{Esribirmos ahora la función de partición para una partícula. Posteriormente vamos a calcularla para todo el sistema}.


Tenemos:


\textbf{La función de partición} de 1 partícula es:


\begin{equation}
\begin{split}
 z &= \sum\limits_{m=-j}^{j} e^{-\beta E_{m}} = \sum\limits_{m=-j}^{j}  e^{ - \beta \left(  -g\mu_{B}H_{m}   \right) }\\
 &= \sum\limits_{m=-j}^{j} e^{\beta g\mu_{\beta} H_{m}} 
\end{split}
\end{equation}

Definimos $x=\beta g\mu_{B}H$


Tenemos:


\begin{equation}
\begin{split}
z &= \sum\limits_{m=-j}^{j} e^{xm} = \underbrace{e^{-xj}e^{xj}}_{1} \sum\limits_{m=-j}^{j}e^{xm}
\end{split}
\end{equation}


Tenemos:


\[ z=e^{-xj} \sum\limits_{m=-j}^{j} e^{x} (m+j) \]

\[  = e^{-xj} \sum\limits_{k=0}^{2j}  e^{xk} \]

Donde $k=m+j$


Recordemos que:


\[   \sum\limits_{k=0}^{m} = \frac{1-r^{m+1}}{1-r}  \]
Tenemos: 


\[   z = e^{-xj}  \left(   \frac{ 1- e^{ x(2j+1)   }   }{1-e^{x}}    \right)  = \frac{ e^{-xj}  -e^{x(j+1)}  }{1- e^{x}}  \]



\[  z= \frac{  - e^{ - \frac{1}{2}x  } }{ -e^{ - \frac{1}{2}x } } \left(  \frac{  e^{-xj} - e^{ x(j+1) } }{ 1 - e^{x} }  \right) \]


\[  z = \frac{ e^{x (j+1/2) }  - e^{-x (j+1/2)}   }{ e^{\frac{x}{2}} - e^{\frac{x}{2}}  } \]

\textbf{Definición y cambio de variable}
\[
z \;\equiv\; \sum_{m=-j}^{j} e^{x m}
\qquad\Longrightarrow\qquad
k = m + j \;\; \Rightarrow \;\; k=0,1,\ldots,2j
\]
\[
z \;=\; e^{-x j}\sum_{k=0}^{2j} e^{x k}
\]

\textbf{Suma geométrica}
\[
\sum_{k=0}^{n} r^{k} \;=\; \frac{1-r^{\,n+1}}{1-r}
\qquad\text{con}\quad r=e^{x},\; n=2j
\]
\[
z \;=\; e^{-x j}\,\frac{1-e^{x(2j+1)}}{\,1-e^{x}\,}
\]

\textbf{Forma simétrica (funciones hiperbólicas)}
Multiplicamos numerador y denominador por \(e^{-x\left(j+\tfrac12\right)}\) para simetrizar:
\[
z \;=\; 
\frac{e^{-x\left(j+\tfrac12\right)}-e^{x\left(j+\tfrac12\right)}}{e^{-x/2}-e^{x/2}}
\;=\;
\frac{-\,2\sinh\!\left(\left(j+\tfrac12\right)x\right)}{-\,2\sinh\!\left(\tfrac{x}{2}\right)}
\;=\;
\boxed{\;z \;=\; \frac{\sinh\!\left(\tfrac{2j+1}{2}\,x\right)}{\sinh\!\left(\tfrac{x}{2}\right)}\;}
\]

\textbf{Comentarios}
\begin{itemize}
  \item La primera línea de tu borrador debe leerse como \(\displaystyle z=\sum_{m=-j}^{j} e^{xm}\) (faltaba el exponente \(xm\))
  \item La identidad final es la forma estándar que aparece en el tratamiento cuántico del paramagnetismo al sumar sobre los niveles \(m=-j,\ldots,j\)
\end{itemize}


Tenemos:



\[  \boxed{ z=   \frac{ \sen h  \left(  \beta \mu_{B} H_{g} \left(  j+1/2 \right)  \right)  }{  \sen h \left( 1/2 \beta \mu_{B}gH  \right)   } } \]


Tenemos:


\begin{figure}[H]
    \centering
    \includegraphics[width=1\textwidth]{2.png}
    \caption{.}
    \label{fig:derivada_total_H}
\end{figure}

\subsubsection{Vamos a graficar la función de Brillouin}


También se puede encontrar que para $y<<$

\[ \langle D_{z} \rangle = N \frac{g^{2} \mu_{B}^{2}  j(j+1) }{ 3K_{B}T } H\]

\begin{figure}[H]
    \centering
    \includegraphics[width=1\textwidth]{3.png}
    \caption{.}
    \label{fig:derivada_total_H}
\end{figure}

\begin{figure}[H]
    \centering
    \includegraphics[width=1\textwidth]{4.png}
    \caption{.}
    \label{fig:derivada_total_H}
\end{figure}


\begin{figure}[H]
    \centering
    \includegraphics[width=1\textwidth]{5.png}
    \caption{.}
    \label{fig:derivada_total_H}
\end{figure}





\end{document}


